% Mechanistic Validity — NEMI 2026 poster, v7
% Five rows by validity type.
% Left: criteria table (bold ID+name, definition, source).
% Right: two case studies per row, padded to match height.
%
%   cd poster && lualatex mechval_poster_v7.tex

\documentclass[25pt, a0paper, portrait]{tikzposter}

\usepackage[T1]{fontenc}
\usepackage{amsmath,amssymb}
\usepackage{array}
\usepackage{booktabs}
\usepackage{enumitem}
\usepackage{microtype}

\tikzposterlatexaffectionproofoff

% ── Colors ────────────────────────────────────────────
\definecolor{darktext}{HTML}{2C3E50}
\definecolor{blockbg}{HTML}{FFFFFF}
\definecolor{constructgreen}{HTML}{27AE60}
\definecolor{measureblue}{HTML}{2980B9}
\definecolor{internalorange}{HTML}{E67E22}
\definecolor{externalpurple}{HTML}{8E44AD}
\definecolor{interpteal}{HTML}{16A085}

% ── Theme ─────────────────────────────────────────────
\usetheme{Default}

\definecolorstyle{MechValColors}{
  \definecolor{colorOne}{HTML}{2C3E50}
  \definecolor{colorTwo}{HTML}{3498DB}
  \definecolor{colorThree}{HTML}{ECF0F1}
}{
  \colorlet{backgroundcolor}{colorThree}
  \colorlet{framecolor}{colorOne}
  \colorlet{titlefgcolor}{white}
  \colorlet{titlebgcolor}{colorOne}
  \colorlet{blocktitlebgcolor}{colorTwo}
  \colorlet{blocktitlefgcolor}{white}
  \colorlet{blockbodybgcolor}{blockbg}
  \colorlet{blockbodyfgcolor}{darktext}
  \colorlet{notefgcolor}{darktext}
  \colorlet{notebgcolor}{colorTwo!12}
  \colorlet{notefrcolor}{colorTwo}
}
\usecolorstyle{MechValColors}

\defineblockstyle{MechValBlock}{
  titlewidthscale=1.0, bodywidthscale=1.0, titleleft,
  titleoffsetx=0pt, titleoffsety=0pt,
  bodyoffsetx=0pt, bodyoffsety=0pt, bodyverticalshift=0pt,
  roundedcorners=12, linewidth=2pt, titleinnersep=10pt, bodyinnersep=12pt
}{
  \draw[rounded corners=\blockroundedcorners, color=blocktitlebgcolor,
        fill=blockbodybgcolor, line width=\blocklinewidth]
    (blockbody.south west) rectangle (blocktitle.north east);
  \ifBlockHasTitle
    \draw[rounded corners=\blockroundedcorners,
          fill=blocktitlebgcolor, color=blocktitlebgcolor,
          line width=\blocklinewidth]
      (blocktitle.south west) rectangle (blocktitle.north east);
  \fi
}
\useblockstyle{MechValBlock}

\definetitlestyle{MechValTitle}{
  width=\paperwidth, roundedcorners=0, linewidth=0pt, innersep=20pt,
  titletotopverticalspace=12mm, titletoblockverticalspace=18mm
}{
  \draw[fill=titlebgcolor, color=titlebgcolor]
    (\titleposleft, \titleposbottom)
    rectangle (\titleposright, \titlepostop);
}
\usetitlestyle{MechValTitle}

\title{%
  \parbox{0.92\linewidth}{\centering\bfseries
    Mechanistic Validity\\[6pt]
    {\large A Validity Framework for Mechanistic Interpretability}}}
\author{Elliot Tower\quad\texttt{elliot@elliottower.ai}}
\institute{}

\begin{document}

\maketitle[width=\paperwidth]

% ════════════════════════════════════════════════════════
% ROW 1: CONSTRUCT VALIDITY
% ════════════════════════════════════════════════════════
\begin{columns}
\column{0.40}

{
\colorlet{blocktitlebgcolor}{constructgreen}
\block{Construct Validity --- Is the target concept well-defined?}{
  A claim is bounded by its weakest validity dimension.
  Five types form a dependency chain: construct validity
  is a ceiling on all others.

  \vspace{4mm}
  \renewcommand{\arraystretch}{1.3}
  {\small
  \begin{tabular}{@{} >{\bfseries}l >{\bfseries}p{0.47\linewidth} >{\bfseries}l @{}}
    \toprule
    \textbf{ID} & \textbf{Description} & \textbf{Source} \\
    \midrule
    C1 Falsifiability & Can the claim be refuted? & Phil.\ of science \\
    C2 Structural plaus. & Physically possible in the architecture? & Phil.\ of science \\
    C3 Convergent & Do independent methods agree? & Measurement theory \\
    C4 Discriminant & Distinguishes from neighboring constructs? & Measurement theory \\
    C5 Nomological & Fits into a broader theory? & Phil.\ of science \\
    C6 Complementation & Are labeled subdivisions functionally distinct? & Genetics \\
    \bottomrule
  \end{tabular}}
}
}

\column{0.30}

{
\colorlet{blocktitlebgcolor}{constructgreen}
\block{C4: Discriminant Validity}{
  Does the measure distinguish this construct from
  neighboring constructs?

  \vspace{4mm}
  \textbf{Knowledge Neurons} (Dai 2022): neurons identified
  by activation magnitude, but the criterion cannot
  separate factual knowledge from linguistic regularity.
  Suppressing these neurons produces off-target effects
  \textbf{worse than suppressing random neurons}.

  \vspace{4mm}
  \textbf{Gender Bias Circuits} (Vig 2020): heads that
  correlate with gendered outputs, but the method never
  tests whether those heads encode bias specifically
  or linguistic competence generally. Bias and competence
  are never separated.

  \vspace{4mm}
  {\footnotesize Source: psychometrics (Campbell \& Fiske 1959).}
}
}

\column{0.30}

{
\colorlet{blocktitlebgcolor}{constructgreen}
\block{C6: Complementation}{
  Are labeled subdivisions functionally distinct?

  \vspace{4mm}
  IOI identifies ``name movers,'' ``S-inhibition heads,''
  and ``backup name movers.'' Whether these labels name
  distinct functional units or overlapping contributions
  to a single mechanism has never been tested.

  \vspace{4mm}
  Benzer's cis-trans test (1955): run two forward passes,
  one with each component intact, and construct a hybrid
  activation (trans configuration). If the hybrid recovers,
  the components serve independent functional roles. If the
  hybrid fails, they participate in the same unit.

  \vspace{4mm}
  {\footnotesize Source: genetics (Benzer 1955, cis-trans test).}
}
}

\end{columns}

% ════════════════════════════════════════════════════════
% ROW 2: MEASUREMENT VALIDITY
% ════════════════════════════════════════════════════════
\begin{columns}
\column{0.40}

{
\colorlet{blocktitlebgcolor}{measureblue}
\block{Measurement Validity --- Does the metric measure what it claims?}{
  \renewcommand{\arraystretch}{1.3}
  {\small
  \begin{tabular}{@{} >{\bfseries}l >{\bfseries}p{0.47\linewidth} >{\bfseries}l @{}}
    \toprule
    \textbf{ID} & \textbf{Description} & \textbf{Source} \\
    \midrule
    M1 Reliability & Repeated measurements agree? & Psychometrics \\
    M2 Baseline sep. & Distinguishable from random baselines? & Psychometrics \\
    M3 Stability & Classification robust to perturbation? & Psychometrics \\
    M4 Calibration & Numbers meaningful on absolute scale? & Psychometrics \\
    M5 Sensitivity & Detects known-true effects? & Psychometrics \\
    M6 Invariance & Consistent across conditions? & Psychometrics \\
    M7 Selection corr. & Multiplicity controlled across $N$ candidates? & Statistics \\
    \bottomrule
  \end{tabular}}

  \vspace{3mm}
  M4 (calibration): attempted in every audited claim,
  confirmed in none.
}
}

\column{0.30}

{
\colorlet{blocktitlebgcolor}{measureblue}
\block{M2: Baseline Separation}{
  Is the score distinguishable from a random baseline?

  \vspace{4mm}
  Three independent failures caught by one criterion:

  \vspace{3mm}
  \textbf{SAE interpretability scores} produce the same
  distribution on random, untrained models as on trained
  ones. The metric cannot distinguish learned features
  from noise.

  \vspace{3mm}
  \textbf{IIA baselines} (Sutter): interchange intervention
  accuracy exceeds majority-class rates before any causal
  structure is identified.

  \vspace{3mm}
  \textbf{SAE reconstruction} (Heap): scores appear strong
  but are indistinguishable from a baseline that ignores
  the learned dictionary entirely.

  \vspace{4mm}
  {\footnotesize Source: psychometrics (floor and ceiling checks).}
}
}

\column{0.30}

{
\colorlet{blocktitlebgcolor}{measureblue}
\block{M6: Measurement Invariance}{
  Does the metric behave consistently across conditions?

  \vspace{4mm}
  \textbf{IOI circuit faithfulness}: the same circuit,
  measured by the same underlying concept (faithfulness),
  produces scores spanning \textbf{below 0\% to over 100\%}
  across six methodological choices.

  \vspace{4mm}
  The headline ``87\% faithfulness'' is not a single number.
  It is a point in a distribution that includes contradictory
  values depending on intervention type, metric formula, and
  ablation target.

  \vspace{4mm}
  If the instrument's reading changes when you swap the
  instrument, the reading is a property of the instrument,
  not the circuit.

  \vspace{4mm}
  {\footnotesize Source: psychometrics (Meredith's measurement
  invariance).}
}
}

\end{columns}

% ════════════════════════════════════════════════════════
% ROW 3: INTERNAL VALIDITY
% ════════════════════════════════════════════════════════
\begin{columns}
\column{0.40}

{
\colorlet{blocktitlebgcolor}{internalorange}
\block{Internal Validity --- Is the causal claim warranted?}{
  \renewcommand{\arraystretch}{1.2}
  {\small
  \begin{tabular}{@{} >{\bfseries}l >{\bfseries}p{0.42\linewidth} >{\bfseries}l @{}}
    \toprule
    \textbf{ID} & \textbf{Description} & \textbf{Source} \\
    \midrule
    I1 Necessity & Is the circuit required? & Neuroscience \\
    I2 Sufficiency & Is the circuit enough? & Neuroscience \\
    I3 Minimality & Every component earns its place? & Causal inf. \\
    I4 Specificity & More effect than matched controls? & Medicine \\
    I5 Rival exclusion & \emph{The} mechanism or \emph{a} mechanism? & Phil.\ of sci. \\
    I6 Double dissoc. & Two interventions, crossed effects? & Neuroscience \\
    I7 Confound ctrl & Alternative explanations ruled out? & Causal inf. \\
    I8 Confound sens. & Unmeasured confounder strength? & Causal inf. \\
    I9 Epistasis & Non-additive component interaction? & Genetics \\
    I10 Rescue & Restoring component recovers behavior? & Genetics \\
    I11 Onset coupling & Mechanism appears with capability? & Med.\ micro. \\
    I12 Offset coupling & Mechanism goes when cap.\ removed? & Med.\ micro. \\
    \bottomrule
  \end{tabular}}
}
}

\column{0.30}

{
\colorlet{blocktitlebgcolor}{internalorange}
\block{I6: Double Dissociation}{
  Do two interventions show crossed effects?

  \vspace{4mm}
  \renewcommand{\arraystretch}{1.3}
  \begin{tabular}{@{}lcc@{}}
    \toprule
    & \textbf{Behavior X} & \textbf{Behavior Y} \\
    \midrule
    Ablate mech.\ A & Large deficit & Small deficit \\
    Ablate mech.\ B & Small deficit & Large deficit \\
    \bottomrule
  \end{tabular}

  \vspace{4mm}
  A selective effect in one row is a single dissociation:
  one intervention affects one behavior. The crossed pattern
  across both rows rules out generalized disruption---neither
  effect can be explained by nonspecific damage.

  \vspace{4mm}
  Across audited claims, only induction heads for token
  copying meet this criterion, achieved by Feucht et al.\
  (2025)---three years after the original paper.

  \vspace{4mm}
  {\footnotesize Source: neuroscience (Shallice 1988).}
}
}

\column{0.30}

{
\colorlet{blocktitlebgcolor}{internalorange}
\block{I9: Epistatic Interaction}{
  Do circuit components interact non-additively?

  \vspace{4mm}
  Ablate A alone, B alone, and both together. Compare the
  joint deficit with the sum of single deficits.

  \vspace{4mm}
  \renewcommand{\arraystretch}{1.3}
  \begin{tabular}{@{}ll@{}}
    \toprule
    \textbf{Pattern} & \textbf{Interpretation} \\
    \midrule
    Joint $=$ sum & Independent contributions \\
    Joint $<$ sum & Redundancy or shared bottleneck \\
    Joint $>$ sum & Synergy or mutual dependence \\
    \bottomrule
  \end{tabular}

  \vspace{4mm}
  Sub-additivity can indicate redundancy \emph{or} a shared
  downstream bottleneck; super-additivity can indicate synergy
  \emph{or} a ceiling effect. The additive null must be
  assessed on an appropriate scale for bounded metrics.

  \vspace{4mm}
  {\footnotesize Source: genetics (Fisher's epistasis).}
}
}

\end{columns}

% ════════════════════════════════════════════════════════
% ROW 4: EXTERNAL VALIDITY
% ════════════════════════════════════════════════════════
\begin{columns}
\column{0.40}

{
\colorlet{blocktitlebgcolor}{externalpurple}
\block{External Validity --- Does the finding generalize?}{
  \renewcommand{\arraystretch}{1.3}
  {\small
  \begin{tabular}{@{} >{\bfseries}l >{\bfseries}p{0.47\linewidth} >{\bfseries}l @{}}
    \toprule
    \textbf{ID} & \textbf{Description} & \textbf{Source} \\
    \midrule
    E1 Intervention reach & Do two different methods agree? & Causal inf. \\
    E2 Prompt general. & Works on diverse prompts? & Causal inf. \\
    E3 Cross-task & Transfers to related tasks? & Causal inf. \\
    E4 Cross-model & Appears in other models? & Causal inf. \\
    E5 Graded response & Partial ablation, partial effect? & Pharmacology \\
    E6 Novel prediction & Predicts untested behaviors? & Phil.\ of sci. \\
    \bottomrule
  \end{tabular}}

  \vspace{3mm}
  E2 (prompt generalization) is disconfirmed for IOI:
  the circuit classification changes across prompt
  distributions.
}
}

\column{0.30}

{
\colorlet{blocktitlebgcolor}{externalpurple}
\block{E1: Intervention Reach}{
  Do two different intervention methods agree?

  \vspace{4mm}
  \textbf{IOI} is the sharpest failure on this criterion.
  Six published evaluations of the same circuit, using
  different intervention methods, produce faithfulness
  scores spanning \textbf{below 0\% to over 100\%}.

  \vspace{4mm}
  When independent methods disagree on the \emph{direction}
  of an effect, the result is a property of the method,
  not the circuit.

  \vspace{4mm}
  The field's most studied mechanism, evaluated by the most
  methods, shows the widest disagreement. Convergent evidence
  from independent methods is what distinguishes the higher
  verdict tiers.

  \vspace{4mm}
  {\footnotesize Source: causal inference (multi-method convergence).}
}
}

\column{0.30}

{
\colorlet{blocktitlebgcolor}{externalpurple}
\block{E5: Graded Response}{
  Does partial intervention produce partial effect?

  \vspace{4mm}
  A dose-response curve between zero and complete ablation
  carries more information than either endpoint alone. The
  curve shape distinguishes threshold effects from graded
  mechanisms, reveals compensation and saturation, and
  detects non-monotonicity invisible at a single dose.

  \vspace{4mm}
  Most published circuit evaluations report a single ablation
  point. The interpolation between intact and fully ablated is
  available in standard activation patching frameworks but is
  almost never recorded or reported.

  \vspace{4mm}
  Genuine mechanisms can produce thresholds; nonspecific
  disruption can produce smooth monotonic decline. The curve
  is informative, not dispositive.

  \vspace{4mm}
  {\footnotesize Source: pharmacology (Hill 1910, dose-response).}
}
}

\end{columns}

% ════════════════════════════════════════════════════════
% ROW 5: INTERPRETIVE VALIDITY + TAKEAWAY
% ════════════════════════════════════════════════════════
\begin{columns}
\column{0.40}

{
\colorlet{blocktitlebgcolor}{interpteal}
\block{Interpretive Validity --- Is the interpretation justified?}{
  \textbf{Novel contribution.} No existing validity framework
  includes these criteria. Interpretive validity asks whether
  the description mode matches the evidence.

  \vspace{4mm}
  \renewcommand{\arraystretch}{1.3}
  {\small
  \begin{tabular}{@{} >{\bfseries}l >{\bfseries}p{0.47\linewidth} >{\bfseries}l @{}}
    \toprule
    \textbf{ID} & \textbf{Description} & \textbf{Source} \\
    \midrule
    V1 Level declaration & At what level is the claim made? & Marr (1982) \\
    V2 Level-evidence & Evidence supports that level? & Messick (1995) \\
    V3 Alternative level & Explained at a different level? & Marr (1982) \\
    V4 Unlicensed labeling & Name imports unmeasured property? & Messick (1995) \\
    V5 Scope declaration & What does the claim not cover? & Reporting std. \\
    \bottomrule
  \end{tabular}}

  \vspace{3mm}
  Activation-level evidence does not license weight-level
  claims. Single-task evidence does not license universal
  claims.
}
}

\column{0.30}

{
\colorlet{blocktitlebgcolor}{interpteal}
\block{V4: Unlicensed Labeling}{
  Does a name import a property that was not measured?

  \vspace{4mm}
  \textbf{Othello ``world model''} (Li 2023): a probe
  recovers board state from internal representations. The
  label ``world model'' imports spatial reasoning, planning,
  and structured representation---none of which were
  contrastively tested. Vafa et al.\ later tested
  counterfactual board states: one checkpoint passes, the
  other fails. The label was premature.

  \vspace{4mm}
  \textbf{Steering ``refusal direction''} (Arditi 2024):
  the direction controls refusal behavior across 13 models.
  But controlling a behavior is evidence for a \emph{lever},
  not evidence for how the model \emph{represents} refusal.
  The label imports representational content the evidence
  does not establish.

  \vspace{4mm}
  {\footnotesize Source: Messick (1995, consequential validity).}
}
}

\column{0.30}

{
\colorlet{blocktitlebgcolor}{colorOne}
\block{Takeaway}{
  {\large\bfseries A mechanistic claim is only as strong
  as its weakest validity dimension.}

  \vspace{4mm}
  \begin{itemize}[leftmargin=*, itemsep=3mm]
    \item \textbf{36} criteria across \textbf{5} validity
      types, drawn from \textbf{8} source disciplines
    \item \textbf{16} claims from \textbf{15} papers audited
    \item The field's most studied circuit (IOI) is
      disconfirmed on measurement stability, measurement
      invariance, and prompt generalization
    \item Same paper, two claims, two tiers: induction heads
      for token copying (Triangulated) vs general in-context
      learning (Disconfirmed)
    \item Interpretive validity (V1--V5) has no counterpart
      in any existing validity framework
  \end{itemize}

  \vspace{4mm}
  \begin{center}
  \texttt{elliot@elliottower.ai}

  \vspace{2mm}
  {\small Paper:}
  \texttt{doi.org/10.5281/zenodo.20478480}
  \end{center}
}
}

\end{columns}

\end{document}
